\documentclass[11pt,a4paper]{article}

\usepackage{amsmath,amssymb,amsthm}
\usepackage{booktabs}
\usepackage{graphicx}
\usepackage{geometry}
\geometry{margin=2.2cm}
\usepackage[dvipsnames]{xcolor}
\usepackage{hyperref}
\hypersetup{colorlinks=true,linkcolor=blue}
\usepackage{fancyhdr}
\usepackage{caption}

\newcommand{\vb}[1]{\mathbf{#1}}
\newcommand{\avg}[1]{\langle #1 \rangle}
\newtheorem{theorem}{Theorem}

\pagestyle{fancy}
\fancyhf{}
\fancyhead[L]{\small Paper Storyline Draft}
\fancyhead[R]{\small\thepage}

\title{\Large\bfseries Time-Domain Local Spectroscopy of Open Dirac Cavities:\\
From Leakage Dynamics to STM-Relevant Spatial Maps}
\author{[Draft --- not for distribution]}
\date{2026-03-28}

\begin{document}
\maketitle

\begin{abstract}
We present a time-domain approach for reconstructing the local density
of states (LDOS) in geometrically open massless Dirac cavities, where
conventional eigenstate-based methods fail due to strongly overlapping
resonances.  By propagating spinor wavepackets in confining potentials
with complex absorbing boundaries, we extract complex resonance
frequencies via matrix-pencil harmonic inversion and reconstruct
STM-relevant spatial maps through multi-propagation LDOS averaging.
We demonstrate that in the overlapping-resonance regime
($\Gamma/\Delta E \sim 1$, enforced by Klein tunnelling), individual
quasi-eigenstate reconstruction is fundamentally ill-posed---not due
to algorithmic limitations, but as a physical consequence of the
energy--time uncertainty relation.  In contrast, the direction-averaged
LDOS remains robust, converging to $C_{2v}$ symmetry errors below
$0.05\%$ with $M = 48$ propagations.  Results are presented for
harmonic and soft-wall ellipses, circles, and Bunimovich stadia,
providing direct theoretical predictions for scanning tunnelling
spectroscopy of graphene quantum dots.
\end{abstract}

\tableofcontents
\vspace{1em}

%======================================================================
\section{Introduction}
\label{sec:intro}
%======================================================================

% --- Step 1: Physical target = local spectral structure ---

The local electronic structure of confined Dirac fermions---how
geometry, confinement potential, and wave dynamics shape the
energy-resolved spatial distribution of states---is a central
question for graphene-based nanostructures.
Scanning tunnelling microscopy (STM) provides direct access to
$dI/dV \propto \text{LDOS}(E, \vb{r})$, making the local density of
states the primary observable linking theory and experiment.
Equally interesting from a fundamental perspective is the nature of
individual quasi-bound states: their spatial wavefunctions, symmetry
properties, and nodal structure carry information about the underlying
classical dynamics (regular vs.\ chaotic) and quantum interference.

% --- Step 2: Bounded systems are well understood ---

For closed or nearly closed billiards---tight-binding models with
hard-wall boundaries, or circular geometries admitting analytical
solutions---computing eigenstates and LDOS is standard practice.
Methods such as kernel polynomial expansion (KPM), recursive
Green's functions, or direct diagonalisation of tight-binding
Hamiltonians have been successfully applied to graphene nanostructures
[refs].  These approaches naturally produce both individual
eigenstates and the LDOS derived from them.

% --- Step 3: Open systems are harder, and eigenstate language
%     becomes less natural ---

The situation changes qualitatively for \emph{open} Dirac cavities
with smooth electrostatic confinement, where carriers can tunnel
through the potential barrier.  For massless Dirac fermions, Klein
tunnelling guarantees unit transmission at normal incidence through
\emph{any} barrier, leading to resonance widths $\Gamma$ bounded
from below by a power law rather than an exponential.  Combined with
the Weyl-law level spacing $\Delta E \propto 1/(EA)$, this gives
\begin{equation}
  \Gamma / \Delta E \sim O(1)
  \quad \text{at all energies, independent of system size,}
\end{equation}
placing the system generically in the \emph{overlapping-resonance
regime}.  In this regime, individual quasi-bound states---while
mathematically well-defined as poles of the scattering matrix---become
increasingly difficult to isolate from finite-time dynamical data,
because neighbouring resonances are temporally indistinguishable
within their own lifetimes.

This does not mean individual states are uninteresting.  On the
contrary, understanding \emph{which} spectral information is reliably
extractable, and which is not, is itself a physically significant
question.

% --- Step 4: What has been done (by capability, not category) ---

Previous theoretical work on relativistic billiards has addressed
eigenstate statistics and scarring in closed systems
[refs], analytical solutions for separable geometries [refs],
and tight-binding simulations of specific nanostructures [refs].
On the numerical spectroscopy side, time-domain wavepacket
propagation combined with harmonic inversion (matrix pencil,
filter diagonalisation) has been used to extract complex resonance
frequencies from autocorrelation signals in molecular and atomic
physics [Mandelshtam \& Taylor 1997; Wall \& Neuhauser 1995],
and more recently in mesoscopic systems [refs].
Dynamic mode decomposition (DMD) and proper orthogonal
decomposition (POD) provide complementary data-driven approaches
that operate on full spatial snapshots rather than scalar signals
[Schmid 2010; refs].
These methods establish that \emph{poles} can be reliably obtained
from short-time dynamics.  However, the question of spatial mode
reconstruction---and its fundamental limitations in the
overlapping-resonance regime characteristic of massless Dirac
systems---has not been systematically addressed.

% --- Step 5: Time-domain as physics, not just technique ---

More broadly, our problem can be viewed as one of resonance
spectroscopy in an open wave system.  In such systems, the
physically relevant spectral structures are not always best
organised in terms of isolated eigenstates, but often appear more
naturally in the language of finite-lifetime response.  Widths are
decay rates, leakage is dynamical, and spectral weight is revealed
through time-dependent response.  This makes the time domain not
merely a convenient numerical route, but a physically natural
framework for open Dirac cavities---conceptually aligned with
emerging ultrafast local spectroscopy techniques [refs].

% --- Step 6: What this paper does ---

In this work, we establish a time-domain workflow for computing
STM-relevant LDOS maps in massless Dirac billiards of arbitrary
geometry, and systematically investigate the limits of individual
mode extraction.  Our main findings are:

\begin{enumerate}
  \item \textbf{Complex resonance frequencies} $(E_n, \Gamma_n)$
        are robustly extracted via matrix-pencil harmonic inversion,
        validated by CAP convergence tests ($<3\%$ variation across
        parameter ranges).

  \item \textbf{Individual mode reconstruction} is possible for
        narrow resonances ($\Gamma/\Delta E \lesssim 0.4$) using
        multiple extraction methods (time filtering, FDM, DMD), but
        progressively fails as $\Gamma/\Delta E \to 1$.
        We quantify this transition and trace it to the energy--time
        uncertainty relation, demonstrating that it is a physical limit
        independent of the extraction algorithm.

  \item \textbf{Multi-propagation LDOS averaging} ($M = 48$
        directions $\times$ positions) produces spatial maps with
        $C_{2v}$ symmetry errors below $0.05\%$ without artificial
        symmetrisation, providing direct theoretical predictions for
        STM experiments.

  \item \textbf{Shape dependence}: results for harmonic and soft-wall
        ellipses, circles, and Bunimovich stadia reveal the interplay
        between integrability, chaos, and Klein tunnelling in
        determining the spatial spectral structure.
\end{enumerate}

%======================================================================
\section{System and numerical method}
\label{sec:method}
%======================================================================

\subsection{2D Dirac equation with confinement}

We consider the 2D massless Dirac Hamiltonian with a scalar confining potential:
\begin{equation}
  H = v_F \begin{pmatrix} 0 & p_x - ip_y \\ p_x + ip_y & 0 \end{pmatrix}
      + V(\vb{r})\,\mathbb{I},
  \label{eq:hamiltonian}
\end{equation}
where $v_F$ is the Fermi velocity and $V(\vb{r})$ provides spatial confinement.
We work in natural units $\hbar = v_F = 1$.

The wavefunction is a two-component spinor $\psi = (\psi_A, \psi_B)^T$ representing
the two sublattice components.  The probability density is
$\rho(\vb{r}) = |\psi_A(\vb{r})|^2 + |\psi_B(\vb{r})|^2$.

\subsection{Confining potentials}

We study three geometries:

\paragraph{Harmonic ellipse.}
$V(\vb{r}) = \tfrac{1}{2}k_\text{harm}\bigl[(x/a)^2 + (y/b)^2\bigr]$
with $k_\text{harm} = 3$, $a = 2$, $b = 1.2$.  The elliptical shape breaks
rotational symmetry while preserving $C_{2v}$.

\paragraph{Soft-wall stadium.}
Following Ref.~[XX], we use the bathtub potential
\begin{equation}
  V(\vb{r}) = \frac{V_0}{1 + \lambda\exp[\mu(1 - f(\vb{r})^2)]},
  \label{eq:softwall}
\end{equation}
where $f^2 = [\max(|x|-a_\text{half},0)^2 + y^2]/b^2$ is the stadium form function,
with $V_0 = 70$, $\lambda = 25$, $\mu = 5$, $a_\text{half} = 1$, $b = 1$.
The interior is nearly flat ($V \approx 0.02$) with smooth walls of width $\sim 1/\mu$.

\paragraph{Harmonic circle.}
$V(\vb{r}) = \tfrac{1}{2}k_\text{harm}(r/r_0)^2$ with $r_0 = 1.55$, chosen to
match the ellipse area ($\pi r_0^2 \approx \pi a b$).

\subsection{Split-operator propagation}

Time evolution uses the split-operator method:
\begin{equation}
  \psi(t + \Delta t) \approx
  e^{-iV\Delta t/2}\, e^{-i T_D \Delta t}\, e^{-iV\Delta t/2}\, \psi(t),
\end{equation}
where $T_D = v_F(\sigma_x p_x + \sigma_y p_y)$ is the kinetic operator, applied
in Fourier space via FFT.  The spatial grid is $256 \times 256$ with $L = 10$,
$\Delta t = 0.01$, $T_\text{max} = 100$.

\subsection{Complex absorbing potential (CAP)}

Open boundaries are modelled by a radial CAP:
\begin{equation}
  W(r) = W_0 \left(\frac{r - r_\text{cap}}{R_\text{max} - r_\text{cap}}\right)^p
  \quad \text{for } r > r_\text{cap},
\end{equation}
with $W_0 = 10$, $r_\text{cap} = 3.5$, $p = 2$.  The propagator includes
$e^{-W\Delta t/2}$ as a real damping factor, ensuring reflection-free absorption
of outgoing waves.

\paragraph{CAP convergence.}
Figure~\ref{fig:cap} shows that extracted resonance widths $\Gamma_n$ are stable
to within 3\% across $W_0 \in [10, 50]$ and $r_\text{cap} \in [2.5, 3.5]$,
confirming that the CAP parameters are in the reflection-free plateau regime.

\begin{figure}[h]
\centering
\includegraphics[width=0.48\textwidth]{../figures/convergence/cap_strength_scan.png}
\includegraphics[width=0.48\textwidth]{../figures/convergence/cap_radius_scan.png}
\caption{CAP convergence: resonance positions $E_n$ and widths $\Gamma_n$ vs.\
CAP strength (left) and CAP radius (right).  Dashed lines mark the default
values.  Individual modes are tracked by nearest energy, not by amplitude ranking.}
\label{fig:cap}
\end{figure}

\subsection{Frequency extraction via matrix pencil}

Complex resonance frequencies $\omega_n = E_n - i\Gamma_n/2$ are extracted from
the scalar autocorrelation function
\begin{equation}
  C_R(t) = \langle\psi_0|\chi_R|\psi(t)\rangle
         = \sum_n |a_n|^2 e^{-i\omega_n t},
\end{equation}
using the matrix-pencil harmonic inversion method.  The analysis mask $\chi_R$
matches the billiard geometry (ellipse, stadium, or circle).

%======================================================================
\section{Local density of states from multi-propagation averaging}
\label{sec:ldos}
%======================================================================

\subsection{LDOS formula}

The local density of states at real energy $E$ is computed as a
Fourier transform of the propagation data:
\begin{equation}
  \text{LDOS}(E, \vb{r}) =
  \left|\sum_k e^{iE t_k}\,\psi_A(\vb{r}, t_k)\right|^2
  + \left|\sum_k e^{iE t_k}\,\psi_B(\vb{r}, t_k)\right|^2.
  \label{eq:ldos}
\end{equation}

This is the modulus squared of a real-frequency Fourier transform.
No mode separation is attempted: overlapping Lorentzians naturally
sum, and the result converges to the physical LDOS.

\subsection{Multi-propagation averaging}

A single propagation from initial state $|\psi_0^\alpha\rangle$ gives
\begin{equation}
  \text{LDOS}_\alpha(E, \vb{r}) =
  \sum_n |\phi_n(\vb{r})|^2\,
  \frac{\Gamma_n/2\pi}{(E-E_n)^2 + (\Gamma_n/2)^2}\,|a_n^\alpha|^2,
\end{equation}
where $|a_n^\alpha|^2 = |\langle\phi_n|\psi_0^\alpha\rangle|^2$ is the
excitation weight for mode $n$.  Different initial states weight
different modes, so averaging over $M$ propagations:
\begin{equation}
  \overline{\text{LDOS}}(E,\vb{r}) = \frac{1}{M}\sum_{\alpha=1}^M
  \text{LDOS}_\alpha(E,\vb{r})
\end{equation}
gives a uniform sampling of all modes.

We use $M = 48$ propagations (8 equally-spaced launch directions
$\times$ 6 positions: origin + ring at $r = 0.4$), with Gaussian
wavepackets of width $\sigma = 0.6$ and peak momentum $|\vb{k}| = 3$.

\subsection{$C_{2v}$ symmetry as a quality indicator}

By the Kramers density identity theorem (proven in [symmetry paper ref]):
each quasi-bound eigenstate density $|\phi_n(\vb{r})|^2$ has exact
$C_{2v}$ symmetry.  Since the LDOS is a positive linear combination of
such densities, it inherits $C_{2v}$ symmetry exactly.

We measure the raw (unsymmetrised) LDOS and find $C_{2v}$ errors
$< 0.05\%$ at all peak energies (Table~\ref{tab:c2v}).  This serves as
a stringent internal consistency check: no artificial symmetrisation is
applied.

\begin{table}[h]
\centering
\caption{$C_{2v}$ symmetry error of raw LDOS average at the 3 strongest
peaks (harmonic ellipse, $M = 48$).}
\label{tab:c2v}
\begin{tabular}{ccc}
\toprule
$E$ & $x$-error & $y$-error \\
\midrule
3.969 & 0.0002 & 0.0000 \\
4.480 & 0.0004 & 0.0000 \\
4.378 & 0.0003 & 0.0000 \\
\bottomrule
\end{tabular}
\end{table}

%======================================================================
\section{Results}
\label{sec:results}
%======================================================================

\subsection{Harmonic ellipse: 2D LDOS maps}

Figure~\ref{fig:ellipse_peaks} shows the LDOS at the 6 strongest
resonance energies.  Clear spatial structure is visible: nodal patterns
reflecting the elliptical confinement, with $C_{2v}$-symmetric density
distributions.

\begin{figure}[h]
\centering
\includegraphics[width=\textwidth]{../figures/ldos/harmonic_ellipse/ldos_peak_maps.png}
\caption{LDOS at the 6 strongest peaks for the harmonic ellipse
($M = 48$, $C_{2v}$ symmetrised).}
\label{fig:ellipse_peaks}
\end{figure}

\subsection{1D line cuts: comparison with STM line scans}

The 2D LDOS can be sliced along symmetry axes to produce
$\text{LDOS}(E, x)$ at $y = 0$ or $\text{LDOS}(E, y)$ at $x = 0$
(Fig.~\ref{fig:linecuts}).  These 2D colour maps directly correspond
to STM line-scan measurements at varying bias voltage.

\begin{figure}[h]
\centering
\includegraphics[width=\textwidth]{../figures/ldos/harmonic_ellipse/ldos_1d_linecuts_300.png}
\caption{1D LDOS line cuts along $y = 0$ (left) and $x = 0$ (right)
for the harmonic ellipse ($n_E = 300$).  Horizontal features correspond
to quasi-bound resonances; their spatial extent reflects the classical
turning curve.}
\label{fig:linecuts}
\end{figure}

\subsection{Shape comparison: ellipse vs circle vs stadium}

Figure~\ref{fig:shapes_spectrum} compares the integrated LDOS spectrum
for three geometries of matched area.  The harmonic ellipse and circle
show similar peak structure (both are integrable classically), while
the soft-wall stadium exhibits broader, less structured features---consistent
with the chaotic classical dynamics of the Bunimovich stadium.

\begin{figure}[h]
\centering
\includegraphics[width=\textwidth]{../figures/ldos/comparison/spectrum_three_shapes.png}
\caption{Integrated LDOS spectrum for three billiard geometries
(harmonic ellipse, harmonic circle, soft-wall stadium), all with
matched area $A \approx 7.5$.}
\label{fig:shapes_spectrum}
\end{figure}

\begin{figure}[h]
\centering
\includegraphics[width=\textwidth]{../figures/ldos/comparison/peaks_three_shapes.png}
\caption{LDOS spatial maps at peak energies for three geometries.
Top: ellipse (elliptical nodal patterns).  Middle: circle (concentric
ring structure).  Bottom: stadium (less regular, reflecting chaotic
dynamics).}
\label{fig:shapes_peaks}
\end{figure}

\begin{figure}[h]
\centering
\includegraphics[width=\textwidth]{../figures/ldos/comparison/linecuts_three_shapes.png}
\caption{1D LDOS line cuts for three geometries.  The ellipse and circle
show clear discrete resonance structure; the stadium shows broader
features with less pronounced gaps.}
\label{fig:shapes_linecuts}
\end{figure}

\subsection{Soft-wall ellipse: realistic confinement}

For comparison with electrostatic quantum dots, we also compute the
LDOS with the soft-wall potential Eq.~(\ref{eq:softwall}) applied to
the ellipse geometry ($a = 2$, $b = 1.2$, same $V_0$, $\lambda$, $\mu$).

\begin{figure}[h]
\centering
\includegraphics[width=\textwidth]{../figures/ldos/comparison/harmonic_vs_softwall.png}
\caption{Harmonic (top) vs soft-wall (bottom) ellipse.  The bathtub
potential produces sharper confinement and slightly different resonance
positions, but the overall spatial structure is similar.}
\label{fig:harmonic_vs_soft}
\end{figure}

\subsection{Stadium at high energy: searching for scars}

Extending the energy range to $E \in [0.5, 25]$ with multiple initial
momenta ($|\vb{k}| = 3, 7, 12, 18$), we search for quantum scarring
in the stadium billiard (Fig.~\ref{fig:scar}).

\begin{figure}[h]
\centering
\includegraphics[width=\textwidth]{../figures/ldos_stadium_scar/scar_candidates.png}
\caption{Scar candidates in the soft-wall stadium at $E > 8$.}
\label{fig:scar}
\end{figure}

\begin{figure}[h]
\centering
\includegraphics[width=\textwidth]{../figures/ldos_stadium_scar/linecuts.png}
\caption{1D LDOS line cuts for the stadium over the extended energy
range $E \in [0.5, 25]$.}
\label{fig:scar_linecuts}
\end{figure}

%======================================================================
\section{Why individual mode extraction fails}
\label{sec:limits}
%======================================================================

\subsection{Klein tunnelling and the $\Gamma/\Delta E$ ratio}

For massless Dirac fermions, Klein tunnelling guarantees unit
transmission at normal incidence through any potential barrier.
This gives a power-law lower bound on the resonance width:
\begin{equation}
  \Gamma \gtrsim \frac{v_F}{k\ell^2},
\end{equation}
which, combined with the Weyl-law level spacing
$\Delta E = 2\pi v_F^2/(EA)$, yields
\begin{equation}
  \frac{\Gamma}{\Delta E} \sim \frac{EL}{2\pi v_F} \sim O(1)
  \quad \text{at the lowest modes.}
\end{equation}

This ratio is \emph{independent} of the billiard size $L$ and
potential shape---a fundamental consequence of the Dirac dispersion
and the absence of exponential evanescent suppression.

\subsection{Fourier resolution limit}

Time-filtered mode reconstruction
$\phi_n(\vb{r}) \propto \sum_k w(t_k) e^{i\omega_n t_k} \chi_R \psi(\vb{r},t_k)$
is equivalent to a windowed Fourier transform with resolution
$\delta E \sim 2\pi\Gamma$ (set by signal lifetime $T_\text{eff} \sim 2/\Gamma$).
When $\Gamma/\Delta E \sim 1$, neighbouring modes overlap within
this resolution window and cannot be separated.

\subsection{Overlap metric: definition and caveats}

To quantify the spatial similarity between an extracted mode
$\phi_\text{method}$ and a reference mode $\phi_\text{ref}$, we use
the normalised spinor overlap:
\begin{equation}
  \mathcal{O} =
  \frac{|\langle\phi_\text{method}|\chi_R|\phi_\text{ref}\rangle|^2}
       {\langle\phi_\text{method}|\chi_R|\phi_\text{method}\rangle\;
        \langle\phi_\text{ref}|\chi_R|\phi_\text{ref}\rangle},
  \label{eq:overlap}
\end{equation}
where
$\langle a|\chi|b\rangle = \sum_{\vb{r}}
[a_A^*\chi\, b_A + a_B^*\chi\, b_B]\,\Delta x^2$.
This quantity equals 1 for identical spatial shapes (up to global
phase and normalisation) and 0 for orthogonal modes.

\paragraph{Important caveat.}
In the massless case, there is no independent ``ground truth'' for
spatial modes.  Our reference $\phi_\text{ref}$ is the time-filtered
mode, which itself is a Lorentzian-weighted superposition of nearby
eigenstates when $\Gamma/\Delta E \sim 1$ (Sec.~\ref{sec:limits}).
Therefore, a high overlap $\mathcal{O} \approx 0.95$ between DMD and
time filtering demonstrates \emph{inter-method consistency}, not
absolute accuracy.  The only intrinsic (reference-free) quality
indicator is $C_{2v}$ symmetry, which is guaranteed by the Kramers
density theorem for exact eigenstates.

\subsection{Multi-channel methods do not remove the barrier}

We systematically tested five mode-extraction methods on the same
propagation data (Table~\ref{tab:methods}; see Appendix~\ref{app:methods}
for method descriptions).

\begin{table}[h]
\centering
\caption{Mode extraction comparison ($m = 0$, harmonic ellipse).
Overlap Eq.~(\ref{eq:overlap}) measured against time-filtered references.
High overlap indicates inter-method consistency, not absolute accuracy
(see text).}
\label{tab:methods}
\begin{tabular}{lccc}
\toprule
Method & Best overlap & Avg overlap & Applicable? \\
\midrule
Time filter & (reference) & (reference) & Baseline \\
FDM & 0.971 & 0.518 & Narrow modes only \\
DMD ($r=30$) & 0.987 & 0.525 & Narrow modes only \\
POD + scalar MP & 0.607 & 0.210 & Poor \\
Block MP (SVD) & 0.349 & --- & Poor \\
\bottomrule
\end{tabular}
\end{table}

All methods give similar results: overlap $> 0.9$ for narrow modes
($\Gamma/\Delta E < 0.4$) and overlap $\approx 0$ for broad modes
($\Gamma/\Delta E > 1$).  The multi-channel advantage (spatial
orthogonality) shifts the threshold slightly but does not remove it.

\subsection{Constrained DMD: pushing extraction to the physical limit}

Since frequencies $(E_n, \Gamma_n)$ are already accurately known from
matrix pencil, one can fix $\lambda_n = e^{-i\omega_n\Delta t}$ in
the DMD framework and solve only for spatial modes (constrained DMD,
or ``Route~C''; see Appendix~\ref{app:methods}).  This combines the
frequency accuracy of the matrix pencil with the full-space spatial
information of the SVD, and outperforms all other methods tested
(Table~\ref{tab:fixedpole}).

\begin{table}[h]
\centering
\caption{Fixed-pole reconstruction ($m = 0$). Route~C achieves
overlap $\approx 1$ on narrow modes.  Tikhonov regularisation and
Laplacian smoothing have no effect (Routes A/B/B+Lap are identical
to bare least squares to 4 digits), confirming the bottleneck is
physical signal content, not numerical conditioning ($\kappa \approx 5$).}
\label{tab:fixedpole}
\begin{tabular}{lcccccc}
\toprule
Method & $E$=3.97 & $E$=4.45 & $E$=3.52 & $E$=2.41 & $E$=3.06 & Avg \\
 & $\Gamma$=0.18 & 0.17 & 0.29 & 0.60 & 0.38 & \\
\midrule
Bare LS / Route A/B  & 0.921 & 0.814 & 0.685 & 0.001 & 0.003 & 0.485 \\
\textbf{Route C} ($r$=40) & \textbf{1.000} & \textbf{1.000} & \textbf{0.718} & 0.002 & 0.004 & \textbf{0.545} \\
DMD ($r$=30)     & 0.987 & 0.921 & 0.716 & 0.002 & 0.006 & 0.527 \\
\bottomrule
\end{tabular}
\end{table}

\subsection{A physical limit on time-domain mode extraction}

The convergence of all algorithms---including Route C---to the same
extractability boundary demands an explanation beyond numerical
limitations.

\paragraph{The argument.}
To separate two modes from time-domain data, one must observe their
beat frequency $|E_n - E_m|$ over at least one full period
$T_\text{beat} = 2\pi / |E_n - E_m|$.  But the signal decays as
$e^{-\Gamma t/2}$, so the signal amplitude at $t = T_\text{beat}$
is suppressed by $e^{-\pi\Gamma/\Delta E}$.  When
$\Gamma/\Delta E \gtrsim 1$, this factor falls below any useful
threshold, and the two modes become indistinguishable from their
time-domain signatures.

This is a limitation of \emph{time-domain observables}, not of the
states themselves.  The quasi-bound states exist as well-defined poles
of the scattering matrix regardless of their overlap.  In principle,
direct diagonalisation of a non-Hermitian Hamiltonian (with explicit
open boundaries) could resolve them.  However, for truly open systems
this requires either:

\begin{itemize}
  \item \emph{CAP-based $H_\text{eff}$}: produces discrete complex
        eigenvalues, but the eigenstates are CAP-parameter-dependent
        and non-orthogonal when resonances overlap.
  \item \emph{Exact open boundaries}: the spectrum is continuous
        (no discrete eigenstates); poles must be found by analytic
        continuation, which is numerically equivalent in difficulty
        to the time-domain approach.
  \item \emph{Green's function at complex energy}: well-defined but
        gives the response function (LDOS), not individual eigenstates.
\end{itemize}

In practice, for the geometries studied here, time-domain propagation
with matrix-pencil pole extraction and constrained DMD mode
reconstruction represents the most practical and general approach.

\paragraph{Contrast with Schr\"odinger billiards.}
For non-relativistic particles, the tunnelling rate decays
\emph{exponentially} with barrier width:
$\Gamma \propto e^{-2\kappa d}$, making $\Gamma/\Delta E$
exponentially small.  Individual mode extraction is then
straightforward.  The qualitative difference---power-law vs.\
exponential suppression of $\Gamma$---is a direct consequence of
the linear Dirac dispersion and the Klein paradox.

\subsection{LDOS bypasses the problem}

The LDOS does not attempt mode separation.  It computes the total
spectral weight at each energy and position, with overlapping
Lorentzians summing naturally:
\begin{equation}
  \text{LDOS}(E, \vb{r}) =
  \sum_n |\phi_n(\vb{r})|^2\,
  \frac{\Gamma_n / 2\pi}{(E - E_n)^2 + (\Gamma_n/2)^2}.
\end{equation}
Each contribution $|\phi_n(\vb{r})|^2$ has exact $C_{2v}$ symmetry
(Kramers density theorem), so the sum inherits it---regardless of
whether individual modes can be resolved.  The LDOS is therefore
a well-defined, robust, and experimentally measurable quantity
even in the strongly overlapping regime where individual mode
extraction fails.

%======================================================================
\section{Discussion}
\label{sec:discussion}
%======================================================================

The results of this work lead to a conceptual distinction that is
worth stating explicitly:

\paragraph{Eigenstate language vs.\ response language.}
For closed or nearly closed quantum systems, the natural theoretical
objects are eigenstates $\phi_n(\vb{r})$, and the LDOS is derived from
them by summation.  For open systems in the overlapping-resonance
regime, this hierarchy inverts: the LDOS (a response function) is the
more fundamental and robust object, while individual eigenstates
become ill-defined as observables---not because they do not exist as
mathematical solutions, but because they cannot be reliably extracted
from any finite-time measurement.

This is not a limitation of our numerical method.  It reflects the
physics of Klein tunnelling: the absence of exponential evanescent
suppression means that resonance lifetimes $\tau_n = 2/\Gamma_n$ are
comparable to the inverse level spacing $1/\Delta E$, so neighbouring
resonances are temporally indistinguishable within their own lifetimes.

\paragraph{Experimental relevance.}
STM measures $dI/dV \propto \text{LDOS}(E, \vb{r})$, which is
precisely the quantity our method computes.  The 2D spatial maps
and 1D line cuts presented here are therefore direct theoretical
predictions for scanning tunnelling spectroscopy of graphene quantum
dots, without requiring the intermediate step of eigenstate
identification.

\paragraph{Adding mass restores mode extractability.}
For systems with a gap (e.g.\ graphene on hBN, $m \neq 0$), the
evanescent barrier restores exponential confinement and
$\Gamma/\Delta E \ll 1$.  In this regime, our pipeline---identical
code, same propagation and extraction---successfully reconstructs
individual modes with overlap $> 0.99$ (Appendix~\ref{app:methods}).
This validates the method and confirms that the failure in the
massless case is physical, not numerical.

\paragraph{Comparison with Schr\"odinger billiards.}
Benchmark calculations for the same geometries with the Schr\"odinger
equation show sharply resolved resonances, extractable eigenstates,
and staircase-like 1D LDOS---the familiar picture of quantum billiards.
The contrast with the massless Dirac case, where peaks are broad and
overlapping, directly illustrates the role of the relativistic
dispersion and Klein tunnelling in reshaping the spectral landscape.

%======================================================================
\section{Conclusion}
\label{sec:conclusion}
%======================================================================

In geometrically open Dirac cavities, single-mode reconstruction
and local spectral reconstruction are fundamentally different in
robustness.  The former fails when $\Gamma/\Delta E \gtrsim 1$
(the generic massless case due to Klein tunnelling); the latter
remains stable and converges with multi-propagation averaging.

We have presented a time-domain workflow---propagation, harmonic
inversion, multi-direction LDOS averaging---that directly produces
STM-relevant spatial maps for massless Dirac billiards of arbitrary
geometry.  The method requires no effective Hamiltonian, no eigenstate
identification, and no assumption about the resonance structure.

For the specific case of graphene quantum dots, the LDOS maps
provide testable predictions: the energy-resolved spatial patterns,
the shape dependence (integrable ellipse vs.\ chaotic stadium), and
the characteristic broadening due to Klein tunnelling.  Comparison
with Schr\"odinger benchmarks quantifies the uniquely relativistic
character of the spectral landscape.

The conceptual takeaway extends beyond the specific system:
\emph{for open quantum systems in the overlapping-resonance regime,
the experimentally natural object is not the eigenstate but the
local spectral response}.

%======================================================================
\appendix
\section{Mode extraction methods}
\label{app:methods}
%======================================================================

We tested five methods for extracting spatial modes from propagation
snapshots $\psi(\vb{r}, t_k)$, $k = 0, \ldots, N_s-1$.  All methods
share the same input data; they differ in how they use frequency
and spatial information.

\subsection{Time filtering (TF)}

The simplest approach: filter at a known complex frequency
$\omega_n = E_n - i\Gamma_n/2$ obtained from matrix pencil:
\begin{equation}
  \phi_n^\text{TF}(\vb{r}) = \sum_k w(t_k)\, e^{+i\omega_n^* t_k}\,
  \chi_R(\vb{r})\,\psi(\vb{r}, t_k),
\end{equation}
where $w(t)$ is a Hann window.  The exponential $e^{+\Gamma t/2}$
compensates signal decay, reweighting late snapshots.

\emph{Limitation}: this is a per-spatial-point Fourier filter with
bandwidth $\sim\Gamma$.  When $\Gamma/\Delta E \sim 1$, neighbouring
modes leak into the filter window.

\subsection{Filter Diagonalisation Method (FDM)}

Builds energy-windowed probes
$\Phi_j(\vb{r}) = \sum_k e^{iE_j t_k}\,\chi_R\,\psi(\vb{r},t_k)$
at $P$ probe energies, then solves a generalised eigenvalue problem:
\begin{equation}
  A\,\vb{c} = z\,B\,\vb{c},\qquad
  B_{jk} = \langle\Phi_j|\Phi_k\rangle\Delta x^2,\quad
  A_{jk} = \langle\Phi_j|\tilde\Phi_k\rangle\Delta x^2,
\end{equation}
where $\tilde\Phi_k$ uses time-shifted snapshots.
SVD regularisation of $B$ (truncate $\sigma < \varepsilon\,\sigma_\text{max}$)
is critical; best parameters: $\varepsilon = 10^{-4}$, $P = 15$,
MP-guided probe energies.

\emph{Advantage}: uses spatial inner products, so modes with
different spatial structure are orthogonal in $B$.
\emph{Result}: frequencies match MP to $<6\%$ ($m=1$); spatial overlap
with TF $>0.99$ ($m=1$), $\sim 0.5$ ($m=0$).

\subsection{Dynamic Mode Decomposition (DMD)}

Operates directly on the full-space snapshot matrix
$X = [\psi(t_0) | \cdots | \psi(t_{N-1})]$ of dimension $D \times N$
($D = 2 n_y n_x = 131\,072$):
\begin{enumerate}
  \item $X_0 = X_{:,0:N-2}$, $X_1 = X_{:,1:N-1}$.
  \item SVD: $X_0 = U\Sigma V^\dagger$, truncate to rank~$r$.
  \item $\tilde{A} = U_r^\dagger X_1 V_r \Sigma_r^{-1}$ \quad ($r\times r$).
  \item Eigendecompose: $\tilde{A}\vb{w} = \lambda\vb{w}$,
        $\omega = i\ln\lambda/\Delta t$.
  \item Exact DMD modes: $\Phi = X_1 V_r \Sigma_r^{-1} W$.
\end{enumerate}

\emph{Advantage}: uses all $D$ spatial dimensions as channels.
The rank $r$ controls a bias--variance tradeoff analogous to
regularisation in machine learning.
\emph{Result}: $m=1$ overlap $>0.999$; $m=0$ similar to FDM
(avg overlap 0.525 at $r=30$).  DMD poles do not form Kramers pairs,
indicating extracted modes are Koopman modes (mixtures of
eigenstates), not individual eigenstates.

\subsection{POD + scalar matrix pencil}

A two-stage hybrid:
\begin{enumerate}
  \item POD (SVD of snapshot matrix) $\to$ $r$ spatial basis
        functions $u_\mu(\vb{r})$ and time coefficients $c_\mu(t)$.
  \item Scalar matrix pencil on each $c_\mu(t)$ independently
        $\to$ frequencies.
  \item Cluster poles across POD components; reconstruct modes
        as $\phi_n = \sum_\mu w_{n\mu}\, u_\mu$.
\end{enumerate}

\emph{Advantage}: decouples spatial compression from spectral
extraction.  \emph{Result}: best overlap 0.61 ($m=0$), worse than
DMD due to pole fragmentation across POD components.

\subsection{Block matrix pencil}

Records $M$ spatial channels
$a_\mu(t) = \langle g_\mu|\psi(t)\rangle$ (spatial patches or SVD
basis), then solves a block Hankel pencil problem for the vector
signal $\vb{a}(t) = \sum_n \vb{v}_n e^{-i\omega_n t}$.

\emph{Advantage}: channel fingerprints $\vb{v}_n$ can separate modes
with the same frequency but different spatial structure.
\emph{Result}: overlap $<0.35$ ($m=0$), limited by the small number
of channels ($M = 19$--$38$) compared to full-space DMD ($D = 131\,072$).

\subsection{Fixed-pole regularised reconstruction}

A fundamentally different approach: frequencies $\omega_n$ are
\emph{fixed} from matrix pencil, and only the spatial modes are
solved for:
\begin{equation}
  \min_{\{\phi_n\}} \sum_k \left\|\psi(\cdot, t_k)
  - \sum_n \phi_n\, e^{-i\omega_n t_k}\right\|^2
  + \lambda\sum_n \mathcal{R}[\phi_n],
\end{equation}
where $\mathcal{R}$ includes Tikhonov regularisation and optional
Laplacian spatial smoothing.  This can be solved in full space
(Route~A) or in a POD-reduced space (Route~B, recommended).

\emph{Advantage}: exploits the strong prior of known poles; reduces
unknowns from $D\times N$ to $r\times N$ (Route~B with POD rank $r$).
\emph{Status}: under evaluation.

\end{document}
